\documentclass{beamer}
\usepackage{graphicx}
\usepackage{xcolor}
\usepackage{tikz}
\usepackage{amsmath}

\definecolor{purpleblue}{RGB}{80, 80, 180}
\definecolor{novelPurple}{RGB}{98, 54, 255}
\definecolor{softPurple}{RGB}{180, 170, 255}
\definecolor{softGray}{RGB}{235, 235, 245}
\definecolor{lightText}{RGB}{90, 90, 110}

\usetheme{Madrid}
\usecolortheme{dolphin}
\graphicspath{{static/course_materials/}{static/hard/}{./}}

\setbeamerfont{title}{size=\Large,series=\bfseries}
\setbeamerfont{frametitle}{size=\LARGE,series=\bfseries}
\setbeamercolor{frametitle}{fg=white, bg=novelPurple}
\setbeamercolor{title}{fg=white, bg=purpleblue}

\title[SAT Hard Question 11]{\textbf{SAT Hard Question\\ Set 11}}
\author{\textbf{Jack Zeng}}
\institute{\textcolor{white}{\textbf{Novel Prep}}}
\date{\today}

\begin{document}

\begin{frame}
    \titlepage
    \vfill
    \begin{flushright}
        \textcolor{purpleblue}{\Huge \textbf{Novel Prep}}
    \end{flushright}
\end{frame}

\begin{frame}
    \begin{tikzpicture}[remember picture, overlay]
        \shade[top color=softPurple, bottom color=softGray]
            (current page.north west) rectangle (current page.south east);
        \node[align=center] at (current page.center) {
            {\Huge\bfseries\textcolor{purpleblue}{Novel Prep}} \\[0.5cm]
            {\Large\itshape\textcolor{lightText}{Phase 2 · Hard Question Practice}}
        };
    \end{tikzpicture}
\end{frame}

\begin{frame}{Overview}
    \tableofcontents
\end{frame}

\begin{frame}
\section{Hard Question Set 11}
    \begin{tikzpicture}[remember picture, overlay]
        \shade[top color=softPurple, bottom color=softGray]
            (current page.north west) rectangle (current page.south east);
        \node[align=center] at (current page.center) {
            {\LARGE\bfseries\textcolor{purpleblue}{Hard Question Set 11}} \\[0.5cm]
            {\Large\itshape\textcolor{lightText}{12 SAT-style challenge problems}}
        };
    \end{tikzpicture}
\end{frame}

% --- Question 1 ---
\begin{frame}{Question 1}
\small
In the given equation, \(a\) and \(b\) are positive constants.
The sum of the solutions to the given equation is \(k(6a + b)\), where \(k\) is constant.
What is the value of \(k\)?
\[
24x^2 - (12a + 2b)x + ab = 0
\]
\vspace{0.35cm}
\begin{enumerate}
    \item[A.] \(\frac{1}{12}\)
    \item[B.] \(-\frac{1}{12}\)
    \item[C.] 12
    \item[D.] -12
\end{enumerate}
\end{frame}
\begin{frame}{Answer 1}
\small
\textbf{Correct Answer:} A
\end{frame}

% --- Question 2 ---
\begin{frame}{Question 2}
\small
The positive number \(a\) is 2,136\% of the sum of the positive numbers \(b\) and \(c\),
and \(b\) is 89\% of \(c\). What percent of \(b\) is \(a\)?
\vspace{0.35cm}
\begin{enumerate}
    \item[A.] 4,536\%
    \item[B.] 2,400\%
    \item[C.] 4,037\%
    \item[D.] 2,225\%
\end{enumerate}
\end{frame}
\begin{frame}{Answer 2}
\small
\textbf{Correct Answer:} A
\end{frame}

% --- Question 3 ---
\begin{frame}{Question 3}
\small
The graphs in the \(xy\)-plane of the following quadratic equations each have
\(x\)-intercepts of \(-2\) and 4. The graph of which equation has its vertex
\textbf{farthest from the \(x\)-axis}?
\vspace{0.35cm}
\begin{enumerate}
    \item[A.] \(y = -7(x + 2)(x - 4)\)
    \item[B.] \(y = \frac{1}{10}(x + 2)(x - 4)\)
    \item[C.] \(y = -\frac{1}{2}(x + 2)(x - 4)\)
    \item[D.] \(y = 5(x + 2)(x - 4)\)
\end{enumerate}
\end{frame}
\begin{frame}{Answer 3}
\small
\textbf{Correct Answer:} A
\end{frame}

% --- Question 4 ---
\begin{frame}{Question 4}
\small
In one hour, the distance, \(d(s)\), in kilometers that a ferry can travel up and down a river
flowing with a constant speed \(s\), in meters per second, is:
\[
d(s) = 10.7 - 1.2s^2
\]
where \(s\) and \(d(s)\) are positive.

Which of the following equivalent expressions for \(d(s)\) contains the speed
of the river in meters per second, as a constant or coefficient, for which
the ferry can travel a distance of 8 kilometers in one hour?
\vspace{0.35cm}
\begin{enumerate}
    \item[A.] \(8 - 1.2(s - 1.5)(s + 1.5)\)
    \item[B.] \(8 + 1.2(2.25 - s^2)\)
    \item[C.] \(8(1 - 0.15s^2) + 2.7\)
    \item[D.] \(11.9 - 1.2(s - 1)^2 - 2.4s\)
\end{enumerate}
\end{frame}
\begin{frame}{Answer 4}
\small
\textbf{Correct Answer:} A
\end{frame}

% --- Question 5 ---
\begin{frame}{Question 5}
\small
A minor league hockey team has been collecting ticket sales data over the past year.
At a current price of \$25 per ticket, an average of 4{,}000 seats are purchased.
For each \$1 increase in ticket price, 100 fewer tickets will be sold.

Which of the following functions best models the total revenue \(y\), in dollars,
from ticket sales when the ticket price increases by \(x\) dollars?
\vspace{0.35cm}
\begin{enumerate}
    \item[A.] \(y = (25 + x)(4000 - 100x)\)
    \item[B.] \(y = (25 - x)(4000 + 100x)\)
    \item[C.] \(y = x(4000 - 100x)\)
    \item[D.] \(y = 4000(25 + x)\)
\end{enumerate}
\end{frame}
\begin{frame}{Answer 5}
\small
\textbf{Correct Answer:} A
\end{frame}

% --- Question 6 ---
\begin{frame}{Question 6}
\small
In the given expression:
\[
34z^{14} + bz^7 + 70,
\]
\(b\) is a positive integer. If \(qz^7 + r\) is a factor of the expression, where \(q\) and \(r\) are positive integers,
what is the greatest possible value of \(b\)?
\end{frame}
\begin{frame}{Answer 6}
\small
\textbf{Final Answer:} $\boxed{2381}$
\end{frame}

% --- Question 7 ---
\begin{frame}{Question 7}
\small
Point \(N\) lies on a unit circle in the \(xy\)-plane and has coordinates \((1, 0)\).
Point \(O\) is the center of the circle and has coordinates \((0, 0)\).
Point \(M\) also lies on the unit circle, and the measure of angle \(\angle NOM\) is \(\frac{266\pi}{3}\) radians.
If the coordinates of point \(M\) are \((a, b)\), where \(a\) and \(b\) are constants, what is the value of \(a\)?
\vspace{0.35cm}
\begin{enumerate}
    \item[A.] \(\frac{1}{2}\)
    \item[B.] \(\frac{\sqrt{3}}{2}\)
    \item[C.] \(-\frac{1}{2}\)
    \item[D.] \(-\frac{\sqrt{3}}{2}\)
\end{enumerate}
\end{frame}
\begin{frame}{Answer 7}
\small
\textbf{Correct Answer:} C
\end{frame}

% --- Question 8 ---
\begin{frame}{Question 8}
\small
For each of two rectangles, the ratio of the rectangle's length to its width is \(10 : 5\).
The width of rectangle \(X\) is 16.5 times the width of rectangle \(Y\).
How does the length of rectangle \(X\) compare to the length of rectangle \(Y\)?
\vspace{0.35cm}
\begin{enumerate}
    \item[A.] The length of rectangle \(X\) is 0.5 times the length of rectangle \(Y\)
    \item[B.] The length of rectangle \(X\) is 2 times the length of rectangle \(Y\)
    \item[C.] The length of rectangle \(X\) is 16.5 times the length of rectangle \(Y\)
    \item[D.] The length of rectangle \(X\) is 33 times the length of rectangle \(Y\)
\end{enumerate}
\end{frame}
\begin{frame}{Answer 8}
\small
\textbf{Correct Answer:} C
\end{frame}

% --- Question 9 ---
\begin{frame}{Question 9}
\small
In the \(xy\)-plane, a circle has center \(C\) with coordinates \((h, k)\).
Points \(A\) and \(B\) lie on the circle. Point \(A\) has coordinates \((h + 1, k + \sqrt{102})\),
and \(\angle ACB\) is a right angle. What is the length of \(\overline{AB}\)?
\vspace{0.35cm}
\begin{enumerate}
    \item[A.] \(\sqrt{206}\)
    \item[B.] \(2\sqrt{102}\)
    \item[C.] \(103\sqrt{2}\)
    \item[D.] \(103\sqrt{3}\)
\end{enumerate}
\end{frame}
\begin{frame}{Answer 9}
\small
\textbf{Correct Answer:} A
\end{frame}

% --- Question 10 ---
\begin{frame}{Question 10}
\small
\[
y = a\left(\frac{6}{b}\right)^{x} - c
\]

How many times does the graph of the given equation in the \(xy\)-plane cross the \(x\)-axis,
where \(a\), \(b\), and \(c\) are positive constants such that \(a > 6\) and \(b > c\)?
\vspace{0.35cm}
\begin{enumerate}
    \item[A.] Zero
    \item[B.] One
    \item[C.] Two
    \item[D.] Three
\end{enumerate}
\end{frame}
\begin{frame}{Answer 10}
\small
\textbf{Correct Answer:} B
\end{frame}

% --- Question 11 ---
\begin{frame}{Question 11}
\small
Two lines intersect at exactly one point, forming two acute angles and two obtuse angles.
The measure of one of these angles is \(9x - 110^\circ\).
Which of the following could \textbf{NOT} be the sum of the measures of any two of these angles?
\vspace{0.35cm}
\begin{enumerate}
    \item[A.] \(-18x + 220^\circ\)
    \item[B.] \(-18x + 580^\circ\)
    \item[C.] \(18x - 220^\circ\)
    \item[D.] \(180^\circ\)
\end{enumerate}
\end{frame}
\begin{frame}{Answer 11}
\small
\textbf{Correct Answer:} A
\end{frame}

% --- Question 12 ---
\begin{frame}{Question 12}
\small
The function \(f\) is defined by:
\[
f(x) = a(3.1^x + 3.1^b),
\]
where \(a\) and \(b\) are integer constants and \(0 < a < b\).
The functions \(g\) and \(h\) are equivalent to function \(f\), where \(k\) and \(m\) are constants.

Which of the following equations displays the \(y\)-coordinate of the \(y\)-intercept of the graph
of \(y = f(x)\) in the \(xy\)-plane as a constant or coefficient?

\begin{itemize}
\item[I.] \(g(x) = a(3.1^x + k)\)
\item[II.] \(h(x) = a(3.1^x) + m\)
\end{itemize}
\vspace{0.35cm}
\begin{enumerate}
    \item[A.] I only
    \item[B.] II only
    \item[C.] I and II only
    \item[D.] Neither I nor II
\end{enumerate}
\end{frame}
\begin{frame}{Answer 12}
\small
\textbf{Correct Answer:} D
\end{frame}

\begin{frame}{}
    \begin{tikzpicture}[remember picture, overlay]
        \shade[top color=novelPurple!40, bottom color=white]
            (current page.north west) rectangle (current page.south east);
        \fill[white, opacity=0.15] (current page.center) circle (5cm);
        \foreach \x in {1,2,3,4,5} {
            \fill[novelPurple, opacity=0.12] (rand*10-5, rand*7-3.5) circle (0.1+\x*0.05);
        }
    \end{tikzpicture}
    \centering
    \vspace{5em}
    {\Huge \textbf{\textcolor{novelPurple}{Thank You!}}} \\[1em]
    {\large \textcolor{gray}{We appreciate your time and focus.}} \\[3em]
    {\LARGE \textbf{\textcolor{novelPurple}{\textsf{Novel Prep}}}}
    \vfill
\end{frame}

\end{document}
